% =====================================================================
\chapter{Логические переключатели}
\label{ch:logic}
% =====================================================================
\chapterintro
  {как создавать «виртуальные» переключатели, которые включаются по условиям: положению
   стика, значению телеметрии, комбинации тумблеров, таймеру}
  {модель с настроенными каналами}

\section{Идея}

\textbf{Логический переключатель} (\emph{logical switch}, \sw{L1}…\sw{L64}) --- это
переключатель, которого нет на корпусе. Его состояние (вкл/выкл) EdgeTX вычисляет по
условию. Логические переключатели можно использовать везде, где можно выбрать обычный
переключатель: в миксерах, входах, таймерах, полётных режимах, специальных функциях
и даже в других логических переключателях.

Примеры задач:
\begin{itemize}
  \item «арм разрешён, только если газ на минимуме»;
  \item «батарея модели ниже 3,5\,В на банку»;
  \item «качество связи упало ниже 50\,\%»;
  \item «кнопка без фиксации работает как тумблер: нажал --- включил, нажал ещё раз --- выключил».
\end{itemize}

\section{Страница Logical Switches}

\mpath{MDL\sep Logical Switches}. У каждой строки:

\begin{center}\small
\renewcommand{\arraystretch}{1.2}
\begin{tabularx}{\linewidth}{@{}p{28mm}X@{}}
\toprule
\menu{Function} & функция сравнения или логики (см. таблицу ниже) \\
\menu{V1}, \menu{V2} & операнды: источники или числа \\
\menu{AND switch} & дополнительное условие: переключатель вкл., только если ещё и этот
  переключатель включён \\
\menu{Duration} & минимальное время, которое переключатель остаётся включённым после
  срабатывания (удлиняет импульс) \\
\menu{Delay} & условие должно выполняться столько секунд, прежде чем переключатель включится
  (фильтр кратких выбросов) \\
\bottomrule
\end{tabularx}
\end{center}

\section{Функции}

\begin{center}\small
\renewcommand{\arraystretch}{1.15}
\begin{tabularx}{\linewidth}{@{}p{24mm}X@{}}
\toprule
Функция & Истинно, когда … \\
\midrule
\code{a=x}, \code{a\textasciitilde x} & значение источника $a$ равно (примерно равно) числу $x$ \\
\code{a>x}, \code{a<x} & источник больше / меньше числа \\
\code{|a|>x}, \code{|a|<x} & модуль источника больше / меньше числа \\
\code{a=b}, \code{a>b}, \code{a<b} & сравнение двух источников \\
\code{Δ≥x}, \code{|Δ|≥x} & источник изменился на $x$ (с момента последнего срабатывания) \\
\code{AND}, \code{OR}, \code{XOR} & логика над двумя переключателями \\
\code{Edge} & переключатель $V1$ был включён на время из диапазона $[t_1; t_2]$ и выключен
  (короткое/длинное нажатие) \\
\code{Sticky} & «защёлка»: включается по $V1$, выключается по $V2$ \\
\code{Timer} & мигает: $V1$ секунд вкл., $V2$ секунд выкл. \\
\bottomrule
\end{tabularx}
\end{center}

\section{Рецепты}

\begin{recipe}{безопасный арм}
Цель: арм (\sw{CH5}) включается, только если в момент перевода \sw{SA} газ был на минимуме.
\begin{enumerate}
  \item \sw{L1}: \code{a<x}, V1 = \sw{Thr}, V2 = $-95$ --- «газ внизу».
  \item \sw{L2}: \code{AND}, V1 = \sw{L1}, V2 = \sw{SA↓} --- «газ внизу и арм включают».
  \item \sw{L3}: \code{Sticky}, V1 = \sw{L2}, V2 = \sw{!SA↓} --- защёлка: включается по \sw{L2},
        выключается, как только \sw{SA} вернули в верхнее положение.
  \item В миксере \sw{CH5}: Source = \sw{MAX}, Weight = $-100$; следующей строкой
        Source = \sw{MAX}, Weight = $+100$, Switch = \sw{L3}, Multiplex = \menu{Replace}.
\end{enumerate}
Если перевести \sw{SA} при поднятом газе, арм не включится; нужно вернуть \sw{SA}, опустить
газ и повторить. После включения арма газ можно поднимать --- защёлка держит состояние.
\end{recipe}

\begin{recipe}{кнопка без фиксации как тумблер}
\sw{SF} на Boxer (или переключатель без фиксации на TX12) --- моментальная кнопка.
\begin{enumerate}
  \item \sw{L5}: \code{Edge}, V1 = \sw{SF↓}, время $[0{,}0; 0{,}5]$\,с --- короткое нажатие.
  \item \sw{L6}: \code{Edge}, V1 = \sw{SF↓}, время $[1{,}0; 3{,}0]$\,с --- длинное нажатие.
  \item \sw{L7}: \code{Sticky}, V1 = \sw{L5}, V2 = \sw{L6}.
\end{enumerate}
Короткое нажатие включает \sw{L7}, длинное --- выключает. Используйте \sw{L7} как обычный
переключатель, например, для подсветки модели.
\end{recipe}

\begin{recipe}{предупреждение о низком напряжении батареи}
\sw{L10}: \code{a<x}, V1 = \sw{RxBt} (напряжение от приёмника/контроллера), V2 = $13{,}6$\,В
(для 4S), \menu{Delay} = 2\,с --- чтобы кратковременные просадки на полном газу не вызывали
ложных тревог. Дальше в специальных функциях: при \sw{L10} --- \menu{Play track}
«батарея разряжена» с повтором каждые 10 секунд (глава~\ref{ch:special}).
\end{recipe}

\begin{recipe}{таймер только в полёте}
\sw{L11}: \code{AND}, V1 = \sw{SA↓} (арм), V2 = \sw{L12}, где \sw{L12}: \code{a>x},
V1 = \sw{Thr}, V2 = $-90$. В таймере \menu{Switch} = \sw{L11}: время идёт, только когда
модель заармлена \emph{и} газ поднят.
\end{recipe}

\section{Отладка}

На главном экране можно пролистать страницу с состоянием всех логических переключателей:
включённые отображаются инверсией. Если условие «не срабатывает», проверьте:
\begin{itemize}
  \item единицы: для телеметрии V2 задаётся в единицах датчика (В, дБм, \%);
  \item знак: у стика газа минимум --- это $-100$, а не 0;
  \item \menu{AND switch} и \menu{Delay} --- не забыты ли там лишние условия.
\end{itemize}

\begin{summary}
\begin{itemize}
  \item \sw{L}-переключатели --- «виртуальные» переключатели по условию.
  \item Сравнения с числом, с другим источником, логика, фронты (\code{Edge}), защёлки
        (\code{Sticky}) и таймеры.
  \item \menu{Delay} фильтрует выбросы, \menu{Duration} удлиняет импульс.
\end{itemize}
\end{summary}
